\chapter{角动量、自旋与统计}

\section{从普通旋转到内禀自旋}

角动量是形式上的主题，但自旋才是真正的目标。自旋是基本粒子最简单也最不直观的性质之一，其结构来自角动量的数学。

角动量是一个矢量。它的方向沿着旋转轴。沿该轴有两个可能的方向，因此需要一种约定来在二者之间做出选择：让右手的四指按旋转方向弯曲，拇指所指的方向就是角动量矢量的方向。这是一个定义，而不是一条定理。

其大小取决于有多少质量存在、该质量如何分布、系统的尺寸以及它的角速度。质量分布和尺寸被归结为转动惯量。因此，一个大物体和一个小物体可以携带相同的角动量，却以非常不同的速率旋转，并具有非常不同的能量。

一个普通的复合物体可以携带两种概念上不同的角动量。一个杯子绕外部轴做圆周运动时具有角动量，因为它的质心在绕该轴运动，即使杯子并没有绕通过自身的任何轴转动。这就是轨道角动量。自旋则意味着绕内部轴的旋转。

干净的定义使用质心系。进入质心静止的参考系，于是总的普通动量消失。在该参考系中仍然残留的任何角动量就是自旋角动量。物体内部不必静止不动；只有它的质心是静止的。

对于一个篮球，这种区分不会造成任何困难。一个静止的篮球和同一个被设定为旋转的篮球，显然被认作处于两种不同状态的同一个物体。量子力学使这幅图景变得复杂，因为角动量是离散的。不存在一系列连续的允许角动量把零角动量的状态与第一个旋转状态连接起来。因此人们可以问：第一个旋转状态是同一个物体在运动，还是一个完全不同的物体。

这种措辞部分是约定性的，但物理问题是精确的：达到第一个非零转动状态需要多少能量？如果间隙很小，受激的物体仍然容易辨认。如果间隙巨大，激发可能会彻底改变该物体。试图给一个原子过多的转动激发，可能会把它的电子甩掉并摧毁这个原子，而不是让同一个原子安静地旋转。

现在考虑一个电子，暂时搁置其已知的内禀自旋。电子能被设定为绕内部轴旋转吗？一个数学上的点没有可见的内部结构可供旋转，但实验并不是通过提供一个可以观察其转动的微小刚性表面来确定电子的大小。如果电子具有极其微小的有限尺寸，那么它的第一个额外转动激发可能只是需要太多能量才能达到。由于能量随尺寸减小而升高，所需的能量可能是天文数字，甚至也许达到普朗克能量的量级。这个估计是刻意试探性的。\footnote{紧接在这个估计之后的一小段文字不够清晰，无法支持进一步的实质性重构。}

因此，对于基本粒子，自旋的大小被当作一个固定的特征。电子并不带有一系列依次增大的内禀自旋大小。如果一个在其他方面类似电子的激发携带不同的总自旋，它就会被归类为一种不同的粒子。

固定的大小并不固定方向。想象一种叫做\emph{自旋子}的粒子，其身份部分地由它角动量的大小来确定。如果它的自旋指向某一根轴，那么旋转不变性要求同一种粒子也可以使其自旋指向另一根轴。没有任何空间方向是内禀优先的。自旋的量保持固定，而它的取向可以改变。于是问题有两个表面上相互冲突的要素：量子化的大小和可自由转动的方向。角动量代数将表明它们如何共存。

\begin{classroomqa}
\lecturetimestamp{00:12:25}
\audiencequestion{为什么让一个小粒子旋转需要能量？}
\lecturerresponse{对于质量为\(M\)、特征半径为\(R\)的物体，转动惯量的量级为
\begin{equation}
I\sim cMR^2,
\end{equation}
其中数值因子\(c\)取决于形状以及质量的分布方式。如果靠近外侧的物质以速率\(v\)运动，其普通动量的量级为\(Mv\)，而角动量等于动量乘以距离：
\begin{equation}
L\sim MvR.
\end{equation}
等价地，\(L=I\omega\)，其中\(v\sim\omega R\)。转动动能为
\begin{equation}
E_{\mathrm{rot}}=\frac{L^2}{2I}.
\end{equation}
在\(M\)和\(L\)固定时，
\begin{equation}
E_{\mathrm{rot}}\sim\frac{L^2}{2cMR^2}\propto\frac{1}{R^2}.
\end{equation}
物体越小，给定角动量所需的能量就越大。这种平方反比依赖关系是经典的。额外的量子力学事实是，角动量以离散的步长变化，因此第一个允许的步长不可能任意小。}
\end{classroomqa}

即将使用的数学是抽象的，起初相当不直观。其意义在于，可由一小组代数规则得出可实验检验的事实。

\section{轨道角动量}

先考虑一个粒子绕原点作圆周轨道运动。其角动量的大小等于动量乘以到原点的垂直距离。在平面内，可以把一种环流方向称为正，把相反方向称为负。这个符号记录了角动量矢量相对于所选轴的方向。

这还不是自旋。该粒子之所以具有角动量，是因为它具有普通动量。现在把两个粒子放在一条直径的两端，让它们以相同的旋转方向运动。它们的普通动量相互抵消，因此质心静止，但它们的角动量指向同一方向并相加。该系统总动量为零而总角动量非零：这就是自旋的经典模型。

设粒子相对于原点的位置为
\begin{equation}
\vec r=(x,y,z)=(x_1,x_2,x_3),
\end{equation}
其动量为
\begin{equation}
\vec p=(p_x,p_y,p_z)=(p_1,p_2,p_3).
\end{equation}
角动量是矢量积
\begin{equation}
\vec L=\vec r\times\vec p.
\end{equation}
交换因子的顺序会使矢量反向：\(\vec p\times\vec r=-\vec r\times\vec p\)。标准约定是\(\vec r\times\vec p\)，其方向由右手定则确定。\footnote{\raggedright 讲座视频帧\nolinkurl{lecture_07_figure_02.png}，00：18：49。\par}

\begin{figure}[h]
\centering
\includegraphics[width=0.72\textwidth]{lecture_07_figure_02.png}
\par
\begin{tikzpicture}[scale=1.0]
\draw[->] (-0.6,0) -- (4.8,0) node[right] {$x$};
\draw[->] (0,-0.6) -- (0,3.8) node[above] {$y$};
\coordinate (P) at (3.2,2.3);
\draw[->,thick] (0,0) -- (P) node[midway,above left] {$\vec r$};
\draw[dashed] (P) -- (3.2,0);
\draw[dashed] (P) -- (0,2.3);
\draw[->,thick] (P) -- ++(-1.0,1.25) node[above] {$\vec p$};
\node at (4.15,2.75) {$\vec L=\vec r\times\vec p$};
\end{tikzpicture}
\caption{对于位于位置\(\vec r\)、动量为\(\vec p\)的粒子，角动量\(\vec L=\vec r\times\vec p\)垂直于它们所在的平面。}
\lectureframe{00:18:49}
\end{figure}

用笛卡尔分量表示，
\begin{equation}
\begin{aligned}
L_x&=yp_z-zp_y,\qquad
L_y=zp_x-xp_z,\\
L_z&=xp_y-yp_x.
\end{aligned}
\end{equation}
只需记住第一个公式并循环置换标签
\[
x\longrightarrow y\longrightarrow z\longrightarrow x.
\]
对分量标签和右侧的标签同时作循环置换，就得到下一个公式，而相对符号不变。这个助记法之所以有用，正是因为符号很容易被意外颠倒。

对于点粒子，\(\vec r\times\vec p\)是相对于所选原点的轨道角动量。对于复合系统，把所有组分的角动量作为矢量相加：
\begin{equation}
\vec L_{\mathrm{tot}}=\sum_a \vec r_a\times\vec p_a.
\end{equation}
当这个和在质心系中计算时，即使
\[
\sum_a\vec p_a=0.
\]
它仍可以保持非零。剩余的角动量就是该复合系统的自旋。

\section{量子算符与正则对易关系}

在量子力学中，可观测量由算符表示。因此，位置、动量以及角动量的每个分量都是算符。这里不需要它们作为微分算符的具体实现。对易关系包含足够的信息来导出角动量谱。

对于两个算符\(A\)和\(B\)，
\begin{equation}
[A,B]=AB-BA.
\end{equation}
如果两个可观测量对易，它们就可以被同时确定。如果它们不对易，量子力学就不允许它们被同时精确确定。

三个位置坐标可以一起测量，因此位置分量对易：
\begin{equation}
[x_i,x_j]=0.
\end{equation}
特别地，每个算符都与自身对易。一个可观测量若与对同一可观测量的另一次测量不相容，那是没有意义的。

动量分量同样对易：
\begin{equation}
[p_i,p_j]=0.
\end{equation}
一个位置分量与不同方向上的动量分量对易，但不与对应的动量分量对易。这三个正则陈述可以合写为
\begin{equation}
[x_i,x_j]=0,\qquad
[p_i,p_j]=0,\qquad
[x_i,p_j]=i\hbar\,\delta_{ij}.
\end{equation}
因此，例如，
\[
[x,p_y]=0,\qquad [x,p_x]=i\hbar,
\]
而交换顺序会改变符号：
\begin{equation}
[p_j,x_i]=-i\hbar\,\delta_{ij}.
\end{equation}

角动量与普朗克常数具有相同的量纲。角动量是距离乘以动量，而不确定性关系具有量纲形式
\begin{equation}
\Delta x\,\Delta p\sim\hbar
\end{equation}
精确不等式中还差一个熟悉的数值因子。因此，角动量将以\(\hbar\)为单位被量子化是很自然的，尽管这种量子化尚未导出。

\section{推导角动量代数}

第一项任务是计算两个角动量分量的对易子。由
\[
L_x=yp_z-zp_y,\qquad
L_y=zp_x-xp_z,
\]
可得
\begin{align}
[L_x,L_y]
&=[yp_z-zp_y,\;zp_x-xp_z] \\
&=[yp_z,zp_x]-[yp_z,xp_z] \notag\\
&\quad{}-[zp_y,zp_x]+[zp_y,xp_z].
\end{align}

\begin{classroomqa}
\lecturetimestamp{00:31:41}
\audiencequestion{能否通过从左到右循环符号、而不是从上到下循环三个方程来记住分量公式？}
\lecturerresponse{可以。两种循环助记法都可以接受。助记法的唯一目的是按正确的循环顺序生成各分量，同时跟踪符号。}
\end{classroomqa}

理解对易子，可以问一个乘积中的因子能否穿过另一个乘积中的因子。在
\[
[yp_z,xp_z],
\]
中，算符 \(y\) 与 \(x\) 和 \(p_z\) 都对易，而 \(p_z\) 与 \(x\) 及其自身都对易。不存在阻碍，因此这一项为零。同样的推理给出
\[
[zp_y,zp_x]=0.
\]
只有两对乘积存活下来。

对于第一对，唯一的阻碍是 \(p_z\) 穿过 \(z\)：
\begin{align}
[yp_z,zp_x]
&=yp_zzp_x-zp_xyp_z \\
&=y(p_zz-zp_z)p_x \\
&=y[p_z,z]p_x \\
&=-i\hbar\,yp_x.
\end{align}
正则对易子的顺序很重要：
\[
[p_z,z]=-[z,p_z]=-i\hbar.
\]

对于另一对存活下来的乘积，\(L_x\) 和 \(L_y\) 的定义中的负号相乘得到正号。唯一的阻碍是 \(z\) 穿过 \(p_z\)：
\begin{align}
[zp_y,xp_z]
&=zp_yxp_z-xp_zzp_y \\
&=x(zp_z-p_zz)p_y \\
&=x[z,p_z]p_y \\
&=i\hbar\,xp_y.
\end{align}
将两项贡献合并，
\begin{align}
[L_x,L_y]
&=-i\hbar\,yp_x+i\hbar\,xp_y \\
&=i\hbar(xp_y-yp_x) \\
&=i\hbar L_z.
\end{align}

这就是重要的简化。两个角动量分量对易产生第三个分量；不再剩下任何位置或动量算符。循环置换给出
\begin{equation}
\begin{aligned}
[L_x,L_y]&=i\hbar L_z,\qquad
[L_y,L_z]=i\hbar L_x,\\
[L_z,L_x]&=i\hbar L_y.
\end{aligned}
\end{equation}
等价地，
\begin{equation}
[L_i,L_j]=i\hbar\,\epsilon_{ijk}L_k.
\end{equation}

代数在角动量算符自身之上封闭。一旦得到这些关系，用来推导它们的 \(x_i\) 和 \(p_i\) 就可以搁置一旁。当总角动量是许多组分的求和时，同样的代数仍然适用。特别地，它适用于质心系，其中剩余的角动量就是自旋。因此，角动量对易子可以抽象地使用，而不必先把自旋构造成微小组成部件的运动。

\(x\)、\(y\) 和 \(z\) 轴的选择是任意的。旋转到任何其他相互垂直的惯性笛卡尔轴，带撇分量满足同样的关系：
\begin{equation}
[L'_x,L'_y]=i\hbar L'_z,
\end{equation}
以及它的循环伙伴。该代数具有旋转不变性，因此物理预言不可能依赖于计算时选择了哪组惯性轴。

\begin{classroomqa}
\lecturetimestamp{00:39:09}
\audiencequestion{自旋角动量是否既独立于坐标轴的平移，也独立于坐标轴的旋转？}
\lecturerresponse{是的。内禀自旋与平移无关。关于某个原点的轨道角动量在原点平移时可能改变，但质心系中剩余的自旋不依赖于该平移。}
\end{classroomqa}

抽象代数可以生成无穷多个恒等式，其中大多数在物理上并无趣味。它的价值出现在某个恒等式成为可测量的事实，或揭示关于测量操作意义的某些东西之时。到目前为止，这一构造仍然主要是代数的。下一步将提取出自旋的离散结构。

\begin{classroomqa}
\lecturetimestamp{00:40:42}
\audiencequestion{如果所选的 \(x\)、\(y\) 和 \(z\) 轴随粒子一起旋转，会怎样？}
\lecturerresponse{讨论仅限于惯性系。旋转坐标系是非惯性的，会引入离心力和科里奥利力之类的效应。角动量在变换到共转系时并非不变：一个在惯性系中明显旋转的系统，在随它一起旋转的参考系中可能显得静止。这里所用的旋转不变性涉及取向不同的惯性笛卡尔轴，而不是随时间旋转的轴。}
\end{classroomqa}

\section{升降算符与量子化轴}

引入两个组合
\begin{equation}
L_+=L_x+iL_y,
\qquad
L_-=L_x-iL_y.
\end{equation}
它们类似于谐振子的产生和湮灭算符。它们的作用是升高和降低一个选定的角动量分量。

\(z\) 轴并无任何特殊之处，但要指定一个分量，就必须选择一根轴。取 \(z\) 为量子化轴，并写出
\begin{equation}
L_z\lvert m\rangle=m\hbar\,\lvert m\rangle.
\end{equation}
因此 \(m\) 是以 \(\hbar\) 为单位的被测 \(z\) 分量。在代数运算过程中，常常方便地令 \(\hbar=1\)；这里将把它明确保留下来。

对易子 \([L_+,L_-]\) 与 \(L_z\) 成正比，但当前的论证并不需要它。首先重要的是每个升降算符与 \(L_z\) 的对易子。由
\[
[L_x,L_z]=-i\hbar L_y,
\qquad
[L_y,L_z]=i\hbar L_x,
\]
可得
\begin{align}
[L_+,L_z]
&=[L_x+iL_y,L_z] \\
&=[L_x,L_z]+i[L_y,L_z] \\
&=-i\hbar L_y+i(i\hbar L_x) \\
&=-\hbar(L_x+iL_y) \\
&=-\hbar L_+.
\end{align}
改变 \(iL_y\) 项的符号，得到
\begin{equation}
[L_-,L_z]=+\hbar L_-.
\end{equation}
等价地，将每个对易子的顺序颠倒，则
\begin{equation}
[L_z,L_+]=\hbar L_+,
\qquad
[L_z,L_-]=-\hbar L_-.
\end{equation}

角动量在任一时刻只能有一个分量是确定的。任意两个不同的分量都不对易，因此不能同时被指定。所选定的确定分量定义了量子化轴；选择 \(z\) 只是为了方便，并非物理上的偏好。

现在用 \(L_+\) 作用于一个 \(L_z\) 本征态。对易子可以改写为
\begin{equation}
L_zL_+=L_+L_z+\hbar L_+.
\end{equation}
因此
\begin{align}
L_z\bigl(L_+\lvert m\rangle\bigr)
&=\bigl(L_+L_z+\hbar L_+\bigr)\lvert m\rangle \\
&=m\hbar L_+\lvert m\rangle+\hbar L_+\lvert m\rangle \\
&=(m+1)\hbar\,L_+\lvert m\rangle.
\end{align}
只要它非零，\(L_+\lvert m\rangle\) 就是一个 \(L_z\) 本征态，其本征值增加了一个 \(\hbar\) 单位。同样地，
\begin{equation}
L_z\bigl(L_-\lvert m\rangle\bigr)
=(m-1)\hbar\,L_-\lvert m\rangle.
\end{equation}
因此 \(L_-\) 将 \(m\) 降低一。

这已经是一个非平凡的结论。无论 \(L_z\) 的谱最终是什么，相邻值之间都相隔整数：
\begin{equation}
\ldots,\ m-2,\ m-1,\ m,\ m+1,\ m+2,\ \ldots
\end{equation}
端点尚未确定。

\section{有限多重态与自旋标记}

一种固定的粒子具有固定的角动量大小。它的角动量可以被重新取向，但无限制地增大其大小最终会产生一个不同的对象。因此，对于该粒子，必定存在一个最大的 \(z\) 分量。

一个带有升算符的阶梯怎么会有顶端？谐振子提供了模型。它的湮灭算符在到达最低态时给出零。在这里，升算符在到达最高态时必须给出零：
\begin{equation}
L_+\lvert m_{\max}\rangle=0.
\end{equation}
\(L_+\) 提升本征值这一说法仅当它产生的矢量非零时才成立。最高权态正是升序列终止的那个例外态。

旋转不变性确定了下端。如果角动量可以最大程度地沿 \(+z\) 方向指向，那么同一个物理对象也可以被旋转，使其最大程度地沿 \(-z\) 方向指向。因此
\begin{equation}
m_{\min}=-m_{\max}.
\end{equation}
谱关于零对称，且间距为单位间隔。存在两个可能的族。

在第一个族中，零出现：
\begin{equation}
\begin{aligned}
m={}&-m_{\max},-m_{\max}+1,\ldots,-1,0,\\
&1,\ldots,m_{\max}-1,m_{\max}.
\end{aligned}
\end{equation}
于是 \(m_{\max}\) 是整数，态数为
\begin{equation}
2m_{\max}+1.
\end{equation}
最大值表征该粒子，称为它的自旋。按惯例记法，
\begin{equation}
s\equiv m_{\max},
\qquad
\lvert s,m\rangle\equiv\lvert m\rangle
\end{equation}
在一个固定自旋多重态内。\footnote{\raggedright 讲座视频帧 \nolinkurl{lecture_07_figure_03.png}，01：01：08。\par}

\begin{figure}[htbp]
\centering
\includegraphics[width=0.70\textwidth]{lecture_07_figure_03.png}
\caption{最高权条件终止了升序列，允许的 \(m\) 值构成一个有限的单位间隔阶梯。}
\lectureframe{01:01:08}
\end{figure}

自旋为零的粒子只有 \(m=0\)。自旋为一的粒子有
\[
m=1,0,-1,
\]
而自旋为二的粒子有五个分量值。希格斯玻色子是自旋为零的一个例子。光子和 \(Z\) 玻色子自旋为一。引力子如果存在，自旋为二。复合对象可以具有自旋三以及高得多的自旋：原子核可以如此，而足够宏观的旋转物体如篮球和星系可以携带任意大的角动量。普通基本粒子列表中缺少的是自旋为三的基本粒子。

另一个对称的、单位间隔的族不包含零。它从
\[
m=\frac12,-\frac12
\]
开始，并以单位步长延伸：
\begin{equation}
\begin{aligned}
m={}&-m_{\max},-m_{\max}+1,\ldots,-\frac32,-\frac12,\\
&\frac12,\frac32,\ldots,m_{\max}.
\end{aligned}
\end{equation}
这里 \(m_{\max}\) 是半整数。最简单的成员自旋为 \(\tfrac12\)，接着是自旋 \(\tfrac32\)、自旋 \(\tfrac52\)，依此类推。因此角动量多重态分为整数自旋族和半整数自旋族。

\begin{classroomqa}
\lecturetimestamp{01:03:46}
\audiencequestion{半整数自旋多重态也有最大值吗？}
\lecturerresponse{有。阶梯同样有一个最高值。对于自旋为 \(\tfrac32\) 的粒子，最大值是 \(m_{\max}=\tfrac32\)；对于自旋为 \(\tfrac52\) 的粒子，最大值是 \(m_{\max}=\tfrac52\)。区别仅在于最大值以及多重态中所有其他值都是半整数。}
\end{classroomqa}

\begin{classroomqa}
\lecturetimestamp{01:04:04}
\audiencequestion{没有自旋为三或更高的基本粒子吗？}
\lecturerresponse{在所讨论的相对论量子框架内，通常列表上的基本粒子自旋为 \(0,\tfrac12,1,\tfrac32,\) 和 \(2\)。复合对象不受此范围限制，可以具有自旋三或大得多的自旋。}
\end{classroomqa}

\begin{classroomqa}
\lecturetimestamp{01:04:50}
\audiencequestion{阶梯上的负值是否代表负自旋？}
\lecturerresponse{不是。负的\(m\)是一个分量值：它表示角动量相对于所选轴指向相反方向。粒子的自旋是非负的最大值\(s=m_{\max}\)。}
\end{classroomqa}

\begin{classroomqa}
\lecturetimestamp{01:05:10}
\audiencequestion{如果分量值是离散的，那么它们之间的所有取向会怎样？}
\lecturerresponse{并非只有两三种可能的空间取向。中间取向由\(m\)基态的量子叠加来表示。下面的自旋\(\tfrac12\)构造明确说明了这一点。}
\end{classroomqa}

\begin{classroomqa}
\lecturetimestamp{01:05:42}
\audiencequestion{那么螺旋运动或圆偏振波呢？}
\lecturerresponse{圆偏振电磁波由圆偏振光子构成，而这些光子携带沿传播轴方向的角动量。``螺旋''一词将旋转与前进运动结合在一起，因此对于被考虑为静止的物体来说，它不是正确的描述。对于运动的粒子，相关的量是螺旋度：其自旋相对于运动方向的取向。}
\end{classroomqa}

\begin{classroomqa}
\lecturetimestamp{01:07:10}
\audiencequestion{是否有已知的基本粒子自旋为\(\tfrac32\)？}
\lecturerresponse{尚未直接确认任何自旋为\(\tfrac32\)的基本粒子。超对称性猜想存在引力子的自旋\(\tfrac32\)伙伴，称为引力微子，但尚未被直接探测到。}
\end{classroomqa}

\begin{classroomqa}
\lecturetimestamp{01:08:01}
\audiencequestion{角动量算符是否已停止作用于原始波函数或右矢空间？}
\lecturerresponse{尚未指定任何原始自旋空间。该代数正被用来确定那个空间。对于固定的自旋\(s\)，它具有由允许值标记的基矢量
\[
\lvert s,m\rangle,\qquad m=-s,-s+1,\ldots,s.
\]
这个内禀自旋空间不同于轨道角动量作为\(\mathbf r\times\mathbf p\)作用的普通空间波函数空间，尽管两者都实现了相同的角动量代数。}
\end{classroomqa}

\section{角动量的大小}

\begin{classroomqa}
\lecturetimestamp{01:09:32}
\audiencequestion{角动量矢量的大小是可观测的吗？}
\lecturerresponse{是的。使用它的平方最为方便，
\[
L^2=L_x^2+L_y^2+L_z^2.
\]
它在最高权态上的值并不简单地是\(s^2\hbar^2\)。即使\(L_z\)是确定的且取最大值，不对易的横向分量也不能同时精确为零，因此它们提供了额外的量子贡献。}
\end{classroomqa}

为了计算该贡献，按升算符先作用于右矢的顺序将阶梯算符相乘：
\begin{align}
L_-L_+
&=(L_x-iL_y)(L_x+iL_y) \\
&=L_x^2+L_y^2+iL_xL_y-iL_yL_x \\
&=L_x^2+L_y^2+i[L_x,L_y] \\
&=L_x^2+L_y^2-\hbar L_z.
\end{align}
最后一步使用了\([L_x,L_y]=i\hbar L_z\)。因此
\begin{equation}
L_x^2+L_y^2=L_-L_++\hbar L_z,
\end{equation}
从而
\begin{equation}
L^2=L_-L_++L_z^2+\hbar L_z.
\end{equation}

将此恒等式应用于最高权态\(\lvert s,s\rangle=\lvert m_{\max}\rangle\)：
\begin{align}
L^2\lvert s,s\rangle
&=\left(L_-L_++L_z^2+\hbar L_z\right)\lvert s,s\rangle \\
&=L_-L_+\lvert s,s\rangle
+s^2\hbar^2\lvert s,s\rangle
+s\hbar^2\lvert s,s\rangle.
\end{align}
第一项消失，因为\(L_+\lvert s,s\rangle=0\)。因此
\begin{equation}
L^2\lvert s,s\rangle
=s(s+1)\hbar^2\lvert s,s\rangle.
\end{equation}
额外的\(+s\)是与不可避免的横向涨落相关的量子修正。经典情况下人们只会预期\(s^2\hbar^2\)。

前几个值使这一区别变得具体：
\begin{align}
s=0 &: \qquad L^2=0,\\
s=1 &: \qquad L^2=2\hbar^2,\\
s=\frac12 &: \qquad L^2=\frac34\hbar^2.
\end{align}
对于\(s=1000\)，无量纲因子为\(1000\times1001\)，在如此大的角动量尺度上与\(1000^2\)差别很小。一般而言，
\begin{equation}
s(s+1)=s^2\left(1+\frac1s\right),
\end{equation}
因此对于非常大的\(s\)，相对量子修正变得可以忽略。经典\(s^2\)结果就是这样出现的。

固定自旋多重态中的每一个\(m\)态都具有相同的\(L^2\)本征值。改变\(m\)改变的是沿所选轴的分量，而不是总大小。这就是在不改变角动量矢量长度的情况下旋转它的量子版本。

\begin{classroomqa}
\lecturetimestamp{01:16:43}
\audiencequestion{\(L^2\)和\(z\)分量能同时已知吗？}
\lecturerresponse{能。总平方与每一个角动量分量都对易：
\[
[L^2,L_i]=0,\qquad i=x,y,z.
\]
因此\(L^2\)和一个选定的分量，例如\(L_z\)，可以一起被指定，尽管两个不同的分量不能。例如，每个电子都有\(L^2=\tfrac34\hbar^2\)，与其取向无关。}
\end{classroomqa}

\begin{classroomqa}
\lecturetimestamp{01:17:41}
\audiencequestion{\(L^2\)是仅对自旋一为零，还是仅对自旋零为零？}
\lecturerresponse{只有自旋零有\(L^2=0\)。自旋一给出
\[
L^2=1(1+1)\hbar^2=2\hbar^2,
\]
而自旋三给出
\[
L^2=3(3+1)\hbar^2=12\hbar^2.
\]}
\end{classroomqa}

\begin{classroomqa}
\lecturetimestamp{01:18:13}
\audiencequestion{为什么多重态中每个\(m\)能级都具有相同的\(L^2\)本征值？}
\lecturerresponse{从最高权态出发并降低它。由于\(L^2\)与每个分量都对易，它也与\(L_-\)对易：
\[
[L^2,L_-]=0.
\]
如果
\[
L^2\lvert s,s\rangle=s(s+1)\hbar^2\lvert s,s\rangle,
\]
那么
\begin{align}
L^2\bigl(L_-\lvert s,s\rangle\bigr)
&=L_-L^2\lvert s,s\rangle \\
&=s(s+1)\hbar^2\,L_-\lvert s,s\rangle.
\end{align}
因此下一个降下来的态具有相同的\(L^2\)本征值。重复这一论证，就会将该本征值一直沿阶梯向下传递。}
\end{classroomqa}

因此，离散的\(m\)值并不意味着只存在离散的一组空间取向。旋转不变性仍然要求存在对应于旋转后方向的态。对于自旋\(\tfrac12\)，这些方向将表现为两个\(L_z\)基态的叠加。

\section{自旋\texorpdfstring{\(\tfrac12\)}{1/2}态与取向}

旋转不变性所要求的不仅仅是离散的\(L_z\)本征态。自旋可以指向任意轴，尽管测量其沿选定轴的分量只给出允许的离散值。\(z\)轴本征态之间的方向由叠加表示。

对于自旋\(\tfrac12\)粒子，只关注自旋自由度。电子还有轨道态，因此它当然并非总共只有两个态；仅其自旋空间就是二维的。选择\(z\)轴基
\begin{equation}
L_z\lvert+\rangle=\frac{\hbar}{2}\lvert+\rangle,
\qquad
L_z\lvert-\rangle=-\frac{\hbar}{2}\lvert-\rangle.
\end{equation}
一般的自旋态为
\begin{equation}
\lvert\psi\rangle
=
\alpha\lvert+\rangle+\beta\lvert-\rangle,
\end{equation}
其中\(\alpha\)和\(\beta\)是复数。测量\(L_z\)得到\(+\hbar/2\)的概率为\(|\alpha|^2\)，得到\(-\hbar/2\)的概率为\(|\beta|^2\)，因此归一化要求
\begin{equation}
|\alpha|^2+|\beta|^2=1.
\end{equation}

整体相位没有物理效应：
\begin{equation}
(\alpha,\beta)\longrightarrow
e^{i\chi}(\alpha,\beta).
\end{equation}
两个复振幅起初有四个实参数。归一化去掉一个，无关的整体相位再去掉一个，剩下两个物理实参数。这恰好是用极角和方位角指定一个方向所需的数目。改变\(\alpha\)和\(\beta\)会旋转自旋取向。

被测分量与其期望值是不同的东西。期望值是系综平均。制备许多处于完全相同状态的电子，对每个电子测量一次\(L_z\)，将结果相加，再除以实验次数。每个单次结果仍然要么是\(+\hbar/2\)，要么是\(-\hbar/2\)，但平均值可以取二者之间的任意值：
\begin{equation}
\langle L_z\rangle
=
\frac{\hbar}{2}\bigl(|\alpha|^2-|\beta|^2\bigr).
\end{equation}

上、下振幅相等给出沿\(x\)轴某一方向的自旋。归一化代表为\footnote{系数比确定了方向。因子\(1/\sqrt2\)提供\( |\alpha|^2+|\beta|^2=1\)所要求的归一化。相对相位\(+i\)和\(-i\)给出相反的\(y\)轴取向；将哪一个符号赋给正\(y\)取决于轴约定。}
\begin{equation}
\lvert x;+\rangle
=
\frac{\lvert+\rangle+\lvert-\rangle}{\sqrt2},
\qquad
\lvert x;-\rangle
=
\frac{\lvert+\rangle-\lvert-\rangle}{\sqrt2}.
\end{equation}
对于这两种态中的任一种，沿\(z\)测量都以相等概率给出两种可能结果，因此
\begin{equation}
\langle L_z\rangle
=
\frac12\left(\frac{\hbar}{2}\right)
+
\frac12\left(-\frac{\hbar}{2}\right)
=0.
\end{equation}
正、负相对符号给出沿\(x\)轴的相反方向。类似地，
\begin{equation}
\frac{\lvert+\rangle+i\lvert-\rangle}{\sqrt2},
\qquad
\frac{\lvert+\rangle-i\lvert-\rangle}{\sqrt2}
\end{equation}
给出两个相反的\(y\)轴取向。\(\alpha\)和\(\beta\)的一般归一化值描述任意轴。说自旋指向某一特定轴，意味着沿该轴测量其分量会给出确定结果。相对于另一条轴，同一个态一般是叠加态。

单次测量结果仍然是量子化的。连续变化并再现经典矢量熟悉几何的是它们的系综平均。

\section{自旋、统计与复合粒子}

质量、电荷和自旋是用于表征粒子的主要性质。自旋还与统计相关。自旋-统计关系是相对论量子力学的一个定理，尽管这里不会证明它：
\begin{equation}
\begin{aligned}
\text{half-integer spin} &\quad\Longleftrightarrow\quad \text{fermion},\\
\text{integer spin} &\quad\Longleftrightarrow\quad \text{boson}.
\end{aligned}
\end{equation}
两个全同费米子不能占据同一个完整量子态。全同玻色子可以占据同一个态，大量玻色子处于一个态时可以像经典波那样集体行为。因此，自旋\(\tfrac12,\tfrac32,\tfrac52,\ldots\)粒子是费米子，而自旋\(0,1,2,\ldots\)粒子是玻色子。

该规则立即推广到复合粒子。两个半整数角动量相加给出整数角动量。更一般地，偶数个费米子组分构成玻色型复合粒子，而奇数个构成费米型复合粒子。加入玻色型组分不会改变这一奇偶规则。

普通氢包含一个自旋为\(\tfrac12\)的质子和一个自旋为\(\tfrac12\)的电子。因此它包含偶数个费米子，整体上是玻色子。氘核包含一个质子和一个中子，两者都是费米子，因此氘核也是玻色子。氘原子在氘核基础上再加一个电子，在此描述层次上给出三个费米子组分，所以氘是费米子。

粒子与反粒子具有相同的质量和自旋，而电荷相反。电子偶素是一个电子和一个正电子相互束缚而成。它包含两个费米子，因此是玻色性的。

\begin{classroomqa}
\lecturetimestamp{01:33:08}
\audiencequestion{玻色子能穿透物体吗？超导性与配对的电子表现得像玻色子有关吗？}
\lecturerresponse{相关的性质并不只是穿透物质的能力；而是玻色子占据同一状态的能力。超流氦给出了更简单的例子。一个普通的氦-\(4\)原子包含两个质子、两个中子和两个电子，即偶数个费米子组分，所以该原子是玻色子。许多氦原子可以占据一个量子态并集体运动。于是超流体可以近似地描绘为由许多原子堆积到同一状态而形成的经典波。这个图像并不精确，但它抓住了这里所需的要点。金属中的电子是费米子，不能与另一个电子共享一个完整状态。然而，电子对是玻色性的复合体，这就是与超导性的联系。这引出一个悖论：如果组成电子不能占据同一状态，那么若干个复合对如何能占据一个状态？解决这个问题需要对复合状态作更仔细的说明，此处暂缓。}
\end{classroomqa}

\begin{classroomqa}
\lecturetimestamp{01:35:46}
\audiencequestion{既然质子和中子是复合粒子，为什么它们具有半整数自旋？}
\lecturerresponse{它们各由奇数个自旋为\(\tfrac12\)的夸克组成。三个费米子组分构成一个费米性复合粒子，所以质子和中子可以具有半整数自旋。}
\end{classroomqa}

\section{泡利不相容原理与原子轨道}

周期表为泡利不相容原理提供了直接的动机。该规则并不是说两个电子永远不能具有相同的轨道运动。而是说两个电子不能占据同一个完整的量子态。

即使在最简单的原子中，这一区别也很重要。从氢出发，增加核电荷并加入第二个电子。这两个电子可以占据同一空间轨道，但它们的自旋态必须不同。相对于所选定的轴，一个可以是自旋向上，另一个是自旋向下。当两个电子占据一个轨道时，它们的自旋组合成自旋单态，因此总自旋相消。一个空间轨道最多只能容纳两个电子，因为只有两个独立的自旋态可供使用。

\begin{classroomqa}
\lecturetimestamp{01:37:53}
\audiencequestion{一个轨道中两个电子的相反自旋可以描绘成啮合的齿轮吗？}
\lecturerresponse{作为一个粗略的图像，可以：啮合的齿轮朝相反方向转动。这只是对相反自旋方向的类比，并不是电子的字面机械模型。}
\end{classroomqa}

``轨道''一词可以在两个不同层次上使用。广义上，轨道运动的一个状态是忽略自旋求解薛定谔方程所得到的状态。它包含的信息比电子仅仅占据哪个壳层更多。

\begin{classroomqa}
\lecturetimestamp{01:38:09}
\audiencequestion{轨道的含义比电子仅仅占据哪个壳层更多吗？}
\lecturerresponse{是的。这里轨道运动指的是电子忽略自旋的薛定谔状态。它包含主量子数\(n\)，后者携带径向信息并大致决定电子到核的距离，还包含轨道角动量量子数。在化学家使用的光谱学记号中，\(S\)、\(P\)、\(D\)和\(F\)分别对应轨道角动量\(l=0,1,2,3\)。}
\end{classroomqa}

因此，主量子数和轨道角动量描述了运动的不同方面。\(S,P,D,F\)标记指的是\(l\)的值，而\(n\)提供单独的径向信息。泡利的规则是：没有两个电子可以具有全部相同的量子数，包括自旋量子数。

\begin{classroomqa}
\lecturetimestamp{01:40:42}
\audiencequestion{量子数能否决定原子中每个壳层可以容纳多少个电子？}
\lecturerresponse{能。它们决定了最大可能的数目，但并不决定某个特定原子在该壳层中实际含有多少个电子。这里不讨论一般的壳层计数问题。}
\end{classroomqa}

对于固定的轨道角动量\(l\)，其分量量子数有
\begin{equation}
2l+1
\end{equation}
个可能取值。这与自旋所用的角动量计数相同。\(S\)态有\(l=0\)，因此有一个轨道态。\(P\)态有\(l=1\)，因此有三个轨道态。\(D\)态有\(l=2\)，因此有五个。

自旋使每个轨道计数加倍，因为每个空间轨道可以在两个相反的自旋态中各容纳一个电子。

\begin{classroomqa}
\lecturetimestamp{01:41:54}
\audiencequestion{\(S\)态容纳两个电子，而三个\(P\)态再容纳六个吗？}
\lecturerresponse{是的。单个\(S\)轨道态最多容纳两个电子。三个\(P\)轨道态各容纳两个，因此再容纳六个。所以，所讨论的\(S\)态和\(P\)态总共容纳\(2+6=8\)个电子。}
\end{classroomqa}

当夸克在束缚系统中运动并相互绕转时，组织原子轨道的同一角动量结构会再次出现，但带有一些差异。
